% !TeX root = ../sections/03_solutions.tex
\usepackage{listing}
\subsection{対策-gnuplot}

\begin{background}
	gnuplot では，まず数学的に求めた式をプログラム上の関数として定義する。
	
	この問題で必要なのは，次の三つである。
	
	\begin{enumerate}
		\item 周期関数 \(f(x)\)
		\item フーリエ係数 \(a_n\)
		\item 部分和 \(S_N(x)\)
	\end{enumerate}
	
	特に，グラフの範囲がデフォルトの \([-10,10]\) のままであるため，
	\(f(x)\) を周期 \(2\pi\) の関数として定義する必要がある。
	
	そのため，任意の \(x\) を区間 \([-\pi,\pi)\) に戻す関数を作る。
\end{background}

\begin{lstlisting}
	# plot.txt の中身
	
	pi = 3.141592653589793
	
	# x を [-pi, pi) に戻す
	p(x) = x - 2*pi*floor((x + pi)/(2*pi))
	
	# 周期関数 f(x)
	f(x) = (abs(p(x)) <= pi/3) ? 0 : \
	(abs(p(x)) <= 2*pi/3) ? pi/2 : pi
	
	# フーリエ係数
	a(k) = -(sin(k*pi/3) + sin(2*k*pi/3))/k
	
	# 第 N 部分和
	S(N,x) = pi/2 + sum [k=1:N] a(k)*cos(k*x)
	
	set yrange [-0.5:4]
	
	plot f(x) title "f(x)" linecolor 1, \
	S(7,x) title "S7(x)" linecolor 2, \
	S(31,x) title "S31(x)" linecolor 3
	
	# graph.png を作るために gnuplot 上で実行する命令
	# これは plot.txt には入れない
	
	set terminal pngcairo size 1000,700
	set output "graph.png"
	load "plot.txt"
	set output
\end{lstlisting}
\begin{strategy}
	gnuplot の問題では，次の順番で考えるとよい。
	
	\begin{enumerate}
		\item まず，数学的にフーリエ係数を求める。
		\item 次に，係数 \(a_n\) を gnuplot の関数として定義する。
		\item 次に，部分和 \(S_N(x)\) を \texttt{sum} を使って定義する。
		\item 最後に，\(f(x)\)，\(S_7(x)\)，\(S_{31}(x)\) を同時にプロットする。
	\end{enumerate}
	
	この問題では，
	\[
	f(x)
	\]
	は階段関数なので，gnuplot では条件分岐
	\[
	\texttt{条件 ? 値1 : 値2}
	\]
	を使って定義する。
	
	また，\([-10,10]\) の範囲で周期関数を描くために，
	\[
	p(x)=x-2\pi\left\lfloor\frac{x+\pi}{2\pi}\right\rfloor
	\]
	を使って，任意の \(x\) を \([-\pi,\pi)\) に戻している。
\end{strategy}
